\subsection{Proposal 4: Error propagation and absorption---the lifecycle of an error}\label{sec:future:lifecycle}

\begin{multicols}{2}
\paragraph{Motivation} Proposals 1--3 treat an error as a point event with a binary consequence: a step is intervened on and the outcome either flips or it does not. The survey data say that what happens \emph{between} the error and the outcome is the common case, not the exception, and that no dataset measures it.
\begin{itemize}[leftmargin=1.1em]
\item Successful runs contain errors. Failure-as-a-Process PASS runs carry 2.3 error events per 22.8 steps, every one of them recovered, with a detour length and a near-miss flag per run (\S\ref{ds:failure_as_process}); in AgentRx the root cause is not the first annotated failure in 27\% of public trajectories and 41\% of Magentic-One logs, because the orchestrator recovers from the earlier ones (\S\ref{ds:agentrx}).
\item Injected errors are usually absorbed. Aegis reports that 23--66\% of its injections leave the outcome unchanged (\S\ref{ds:aegis}); a reviewing Inspector agent recovers up to 96.4\% of injected faulty-agent errors \cite{huang2024resilience}; claim-level hallucinations attenuate rather than amplify along a three-agent chain (0.422 $\to$ 0.272) \cite{jamshidi2026cascade}. All four injected datasets of Section~\ref{sec:g2} discard the absorbed branch.
\item Recovery is fast or not at all. The fix window $t_{\mathrm{lock}}-t_{\mathrm{err}}$ has median 1 step (239 of 1,184 runs above 10); of the failed runs, 58.7\% never notice the error and 34.9\% notice it and apply the wrong fix; 28.0\% fail silently, with no externally observable signal at all.
\item Recovery is trainable. Recovery-targeted fine-tuning raises a computer-use agent from 23.6 to 35.2 at matched tokens \cite{cuarecovery2026}, and a perturbation-recovery rate falls by $\approx$37\% under implicit tool failures while task success moves much less \cite{zhu2026toolmaze}.
\end{itemize}
The hypothesis this proposal makes measurable is that a robust agent system is one in which an error introduced at step $k$ is detected or corrected by a later step with high probability and short delay, so that the trajectory still succeeds. Two fields have measured exactly this for fifty years and supply the vocabulary. In dependable computing \cite{avizienis2004dependability} a \emph{fault} (the perturbation) produces an \emph{error} (a wrong internal state) which may or may not become a \emph{failure} (a wrong output); the error can stay \emph{dormant}, be \emph{masked}, be \emph{detected}, be \emph{recovered} or \emph{propagate}; fault injection measures the \emph{coverage} (the probability that an activated fault is handled) and the \emph{error latency} (the delay until it is) \cite{arlat1990fault,powell1995coverage}, and the PIE model factorises failure probability into execution, infection and propagation \cite{voas1992pie}. In ecology and infrastructure engineering, \emph{resilience} is the capacity to return to the basin of attraction after a perturbation \cite{holling1973resilience}, measured by the depth and duration of the performance dip \cite{bruneau2003resilience}. The programme below transfers these quantities to agent trajectories with the notation of Section~\ref{sec:future:prelim}, and gives the quantity $P_{\mathrm{detect}}(k{+}i\mid k)$ of the hypothesis a precise home.

\paragraph{Closest prior work} Besides the three recovery studies above \cite{huang2024resilience,jamshidi2026cascade,zhu2026toolmaze}, four strands contain parts of the programme. \emph{Error-propagation profiling} of software modules \cite{hiller2004epic} measures, per module, how often a wrong input produces a wrong output (permeability) and how many consumers it reaches (exposure); EMFA \cite{liao2026emfa} applies an error-propagation model to Who\&When but fits it to failed runs only. \emph{Recoverability in safe RL} \cite{moldovan2012safe,grinsztajn2021reversibility} defines a state as recoverable if the agent can return to a safe set, and learns reversibility from trajectories; the point of no return of task T3 is that notion on agent runs. \emph{Multi-state event-history analysis} \cite{putter2007multistate,andersen1982cox,aalen1978empirical} estimates the competing exits of a subject from a state under censoring, which is what an error that is still live when the run ends requires. \emph{Quickest change detection} \cite{page1954cusum,lorden1971,tartakovsky2014sequential} gives the detection-delay versus false-alarm trade-off for the online monitor of task U3. None of them has been applied to an agent trajectory with a replayable environment, which is what makes every quantity below measurable rather than judged.

\paragraph{Three event times, not one} Fix a successful trajectory $\tau$, a step $k$ and an injection $\delta\in\Delta_{\mathrm{inj}}$ (Table~\ref{tab:vocab}), and let $\omega\sim\mathrm{Roll}(\tau_{<k}\oplus\delta(s_k))$ be one continuation. Three things can happen to the injected error, at three different lags $i\ge1$ after $k$:
\begin{itemize}[leftmargin=1.1em]
\item \textbf{detection} $D$: the first lag at which a step of $\omega$ \emph{refers} to the inconsistency (an observation contradicts the perturbed content, a verifier flags it, the agent says so)---the only one of the three that is visible in the text and the quantity $P_{\mathrm{detect}}(k{+}i\mid k)=\Pr[D{=}i]$;
\item \textbf{recovery} $R$: the first lag after which the perturbed content no longer governs the run---operationally, deleting $\delta(s_k)$ from the context of every step after $k{+}R$ leaves the continuation distribution unchanged, and the success probability of the prefix has returned to its unperturbed level;
\item \textbf{lock-in} $L$: the first lag from which no single-step repair in $\Delta_{\mathrm{fix}}$ rescues the run---Definition~\ref{def:fa} applied to $\omega$.
\end{itemize}
$R$ and $L$ compete: at most one of them occurs, and a run can also end with the error live but harmless (success without recovery), which is a \emph{censored} observation of both. $D$ is not an exit; a detected error can still lock in (Failure-as-a-Process's ``aware, wrong fix'', 34.9\%), and an error can be recovered without ever being detected (a second retrieval returns the right value and the wrong one is forgotten). Failure-as-a-Process's $t_{\mathrm{obs}}$, $t_{\mathrm{lock}}$ and the detour length of its PASS runs are the natural-failure versions of $D$, $L$ and $R$; here they are measured on injected errors where $k$ is known by construction.

\begin{definition}[Error fate]\label{def:fate}
The \emph{fate} of the injected error in $\omega$ is one of: \emph{dormant} (no later step consumes $\delta(s_k)$: deleting it from every later context changes nothing); \emph{masked} (consumed, but the next decision is the same as under $s_k$); \emph{corrected} ($D<\infty$ and $R<\infty$); \emph{overwritten} ($R<\infty$, $D=\infty$); \emph{tolerated} (consumed, never recovered, $S{=}1$); \emph{fatal} ($L<\infty$, $S{=}0$). The \emph{coverage} of the system at $(\tau,k,\delta)$ is $C(\tau,k,\delta)=\Pr[\text{fate}\neq\text{fatal}]=p_\delta(\tau,k)$ of Definition~\ref{def:effect}; the fate distribution is its decomposition.
\end{definition}

Proposal~1's injection question (Q2) and its fragility curve $r(k/T)$ report only $C$, i.e.\ the product of the three PIE factors. Two systems with the same $r$ can differ entirely in how they get it: one by dormancy (the task can be solved without the corrupted content), the other by correction. The fate distribution separates the two, and only the second is robustness in the sense of the hypothesis.

\begin{definition}[Error lifetime]\label{def:lifetime}
Let $Z_i\in\{\text{live},\text{recovered},\text{locked in}\}$ be the state of the error at lag $i$, with $Z_0=\text{live}$. The \emph{cause-specific hazards}
\begin{align*}
h_{\mathrm{rec}}(i\mid\tau,k,\delta)&=\Pr[R{=}i\mid Z_{i-1}{=}\text{live}],\\
h_{\mathrm{lock}}(i\mid\tau,k,\delta)&=\Pr[L{=}i\mid Z_{i-1}{=}\text{live}]
\end{align*}
and the survival of the live state $\bar S^{\mathrm{err}}_i=\prod_{j\le i}(1-h_{\mathrm{rec}}(j)-h_{\mathrm{lock}}(j))$ give the \emph{cumulative incidences} $F_{\mathrm{rec}}(i)=\sum_{j\le i}\bar S^{\mathrm{err}}_{j-1}h_{\mathrm{rec}}(j)$ and $F_{\mathrm{lock}}(i)$ likewise; $F_{\mathrm{rec}}(\infty)$ is the recovery probability and $\mathbb{E}[R\mid R<\infty]$ the mean time to recovery. The same construction with $D$ in place of $R$ gives the detection-latency distribution.
\end{definition}

The error of Definition~\ref{def:lifetime} is the \emph{subject} of a survival analysis, which the step of Definition~\ref{def:hazard} is not: Proposal~2's model is the special case $h_{\mathrm{rec}}\equiv0$, in which every live error is eventually fatal (Assumption~\ref{as:absorb}). The multi-state model in Fig.~\ref{fig:lifecycle} contains both, and the generalisation of the hazard model that fits PASS runs with recovered errors is the one with a return transition.

\paragraph{Three derived quantities}
\begin{itemize}[leftmargin=1.1em]
\item \emph{Resilience triangle.} Let $V^0_t=p_0(\tau,t{+}1)$ be the success probability of the unperturbed prefix (the survival function $\bar S_t$ of Proposal~2) and $V^\delta_t$ the same for $\omega$. The \emph{resilience loss} $\mathcal{R}(\tau,k,\delta)=\sum_{i\ge0}\big[V^0_{k+i}-V^\delta_{k+i}\big]$ is the area between the two curves \cite{bruneau2003resilience}; its depth $\max_i(V^0_{k+i}-V^\delta_{k+i})$ and its rapidity (the first $i$ at which the gap closes to $\alpha$) are the robustness and recovery speed of the system, as two numbers rather than one. Recovery shows as a dip and a rise of $V^\delta$; lock-in as a dip without a rise.
\item \emph{Per-agent profile.} For an agent $a\in\mathcal{A}$ of a multi-agent system, following \cite{hiller2004epic}, four rates---\emph{generation}, \emph{permeability}, \emph{containment} and \emph{exposure}:
\begin{align*}
g_a&=\Pr[u_t\ \text{wrong}\mid a_t{=}a,\ \text{clean input}],\\
\pi_a&=\Pr[u_t\ \text{wrong}\mid a_t{=}a,\ \text{live error in input}],\\
c_a&=\Pr[\text{flags or fixes it}\mid a_t{=}a,\ \text{live error in input}],\\
x_a&=\mathbb{E}[\#\,\text{later steps consuming }u_t\mid a_t{=}a].
\end{align*}
``How prone is each agent to committing to errors'' is $\pi_a$; ``how prone to committing errors'' is $g_a$; a robust team is one with low $\pi$ and high $c$ on the agents with high $x$. The matrix $M_{ab}=\pi_b\cdot\mathbb{E}[\#\,b\text{-steps consuming one }a\text{-output}]$ is the mean-offspring matrix of a multi-type branching process \cite{harris1963branching}: $\rho(M)<1$ means a single wrong message dies out with probability one, and the extinction probability solves the usual fixed-point equation, so absorption can be \emph{predicted} from pairwise measurements and compared with the empirical $C$. The topology result of \citet{huang2024resilience} (hierarchical 5.5\% loss versus 10.5--23.7\% for flat) is a statement about $\rho(M)$.
\item \emph{Checkpoint model.} Let $\mathcal{C}(\tau)$ be the steps whose role is verification (verifier-agent turns, test runs, re-retrievals, orchestrator ledger updates), each with coverage $c_j$. The null model $C(\tau,k,\delta)\approx1-\prod_{j\in\mathcal{C},\,j>k}(1-c_j)$ predicts that $r(k/T)$ is a step function aligned with the checkpoints rather than a smooth decay, that removing a verifier moves it by a measurable $c_j$, and that the late-position skew of the injected datasets (Fig.~\ref{fig:errpos}) is the regime with no checkpoint left.
\end{itemize}
\end{multicols}

\begin{figure}[h]\centering
\begin{tikzpicture}[x=1cm,y=1cm,font=\scriptsize,>=Stealth,
  st/.style={draw=rule!80!black,fill=surface,rounded corners=2pt,minimum height=0.5cm,minimum width=1.55cm,inner sep=2pt,align=center},
  bad/.style={draw=critical,fill=critical!8,rounded corners=2pt,minimum height=0.5cm,minimum width=1.55cm,inner sep=2pt,align=center},
  ok/.style={draw=good!70!black,fill=good!8,rounded corners=2pt,minimum height=0.5cm,minimum width=1.55cm,inner sep=2pt,align=center}]
% multi-state diagram
\node[st] (track) at (0,0) {on track};
\node[bad] (live) at (3.0,0) {error live\\(born at $k$)};
\node[ok] (rec) at (6.2,0.95) {recovered};
\node[bad] (lock) at (6.2,-0.95) {locked in};
\node[ok,draw=good!50,fill=good!4] (succ) at (9.2,0.95) {$S=1$};
\node[bad,draw=critical!50,fill=critical!4] (fail) at (9.2,-0.95) {$S=0$};
\draw[->,ink2] (track) -- node[above,text=ink2]{$g_a$ or $\delta\in\Delta_{\mathrm{inj}}$} (live);
\draw[->,good!70!black] (live) -- node[above left=-1pt,text=good!60!black]{$h_{\mathrm{rec}}(i)$} (rec);
\draw[->,critical] (live) -- node[below left=-1pt,text=critical]{$h_{\mathrm{lock}}(i)$} (lock);
\draw[->,ink2] (rec) -- node[above,text=ink2]{no further error} (succ);
\draw[->,ink2] (lock) -- (fail);
\draw[->,ink2,dashed] (rec.north) .. controls (6.2,2.0) and (0,2.0) .. node[above,text=ink2,pos=0.5]{next error (recurrent)} (track.north);
\draw[->,ink2,dotted] (live.east) .. controls (4.6,0.3) and (7.4,0.4) .. node[below,text=ink2,pos=0.75]{tolerated / dormant (censored)} (succ.west);
\node[text=ink2,align=center,anchor=north] at (2.3,-0.95) {detection $D$: a text-visible\\event inside the live state};
% value curve with resilience triangle
\begin{scope}[shift={(11.6,-1.35)}]
\draw[->,rule!80!black] (0,0) -- (4.3,0) node[right,text=ink2]{$t$};
\draw[->,rule!80!black] (0,0) -- (0,2.0) node[above,text=ink2]{$V_t$};
\draw[s1!80,line width=0.7pt] (0,1.6) -- (4.0,1.6);
\draw[critical!80,line width=0.7pt] (0,1.6) -- (1.0,1.6) -- (1.5,0.55) -- (2.4,0.7) -- (3.1,1.5) -- (4.0,1.6);
\draw[critical!50,line width=0.7pt,dashed] (2.4,0.7) -- (3.1,0.3) -- (4.0,0.25);
\fill[critical!15,opacity=0.6] (1.0,1.6) -- (1.5,0.55) -- (2.4,0.7) -- (3.1,1.5) -- cycle;
\draw[ink3,dashed] (1.0,0) -- (1.0,1.6) node[below,at start,text=ink2]{$k$};
\draw[ink3,dashed] (3.1,0) -- (3.1,1.5) node[below,at start,text=ink2]{$k{+}R$};
\draw[<->,ink2] (1.5,0.55) -- (1.5,1.6) node[midway,right,text=ink2]{depth};
\node[text=ink2] at (2.1,1.78) {rapidity};
\node[text=critical!70!black,anchor=west] at (3.3,0.45) {lock-in: no rise};
\node[text=s1!80!black,anchor=west] at (0.1,1.85) {$V^0$};
\end{scope}
\end{tikzpicture}
\caption{Left: the multi-state model of an error's lifetime. The live state is exited by recovery or lock-in with competing hazards; a run that ends with the error live but harmless is a censored observation. Proposal~2's hazard model is the special case without the recovery transition. Right: the same on the success-probability curve $V_t$ of the prefix (Proposal~2's $\bar S_t$): recovery is a dip followed by a rise, lock-in a dip without one, and the shaded area is the resilience loss $\mathcal{R}$ with its depth and rapidity.}
\label{fig:lifecycle}
\end{figure}

\begin{multicols}{2}
\paragraph{Research questions}
\begin{enumerate}[leftmargin=1.4em,itemsep=1pt]
\item \textbf{Lifetime.} What are $F_{\mathrm{rec}}$, $F_{\mathrm{lock}}$ and the detection-latency distribution as functions of lag, and how do they depend on position $k/T$, error mode, fault side and system? Is the one-step median fix window of natural failures also the recovery time of injected errors, or is it a selection effect of failed runs?
\item \textbf{Mechanism.} What fraction of absorption is dormancy, masking, correction and overwriting? Only the last two are the hypothesised capability; the first two are task slack.
\item \textbf{Checkpoints.} Is $r(k/T)$ a step function aligned with $\mathcal{C}(\tau)$, and does removing or adding a verifier move it by the predicted $c_j$?
\item \textbf{Detected or overwritten.} How often is $R<\infty$ with $D=\infty$, and does a detected-but-unfixed error ($D<L<\infty$) look different in the text from a never-detected one?
\item \textbf{Per-agent profile.} What are $(g_a,\pi_a,c_a,x_a)$ for the roles of the surveyed systems, is $\rho(M)<1$, and does the branching-process absorption probability predict the empirical $C$?
\item \textbf{Dose--response.} With a severity-graded injector $\delta_\sigma$ (size of a numeric error, fraction of a page truncated, number of corrupted arguments), what is $\Pr[\text{fatal}\mid\sigma]$, and is there a critical severity $\sigma^\star(k)$, the basin size of the resilience picture \cite{holling1973resilience}?
\item \textbf{Epistasis.} For two injections at $(k_1,k_2)$, is $\log p_{\delta_1\delta_2}-\log p_{\delta_1}-\log p_{\delta_2}+\log p_0$ zero (independent), negative (compounding) or positive (the first error triggers a check that catches the second)? Minimal repair families \cite{li2026mrfr} show the fix side of this is non-additive.
\item \textbf{Cost.} What is the resilience loss $\mathcal{R}$, in steps and tokens, of a recovered error, and is recovering cheaper than restarting from $\tau_{<k}$?
\item \textbf{Identifiability from the log.} Can a recovered error be found without rollouts, from the dip-and-rise signature of a learned $V_t$ or from a hidden state that switches back, and with what delay can an online monitor tell a transient dip from a lock-in?
\end{enumerate}

\paragraph{Measurement tasks}
\begin{description}[leftmargin=1.4em,style=nextline,itemsep=1pt]
\item[E1 --- Fate screen.] The injection half of Proposal~1's screen with every continuation \emph{kept} and labelled with Definition~\ref{def:fate}. The consumption test is a second, cheap counterfactual: for each step $k{+}i$ of $\omega$, re-issue that single LLM call with $\delta(s_k)$ removed from its context and compare the output (one call per step, not a rollout). Detection $D$ is a judge pass over the text with the injection known; recovery $R$ is the first lag at which the consumption test stays negative for the rest of the run and $V^\delta$ has returned to $V^0-\alpha$; lock-in $L$ is estimated by the surrogate $f$ of task T1 and verified by rollouts on a subset. Output: the fate distribution per system, mode, position and fault side, and the triples $(D,R,L)$ with censoring indicators.
\item[E2 --- Lifetime curves.] On the E1 output, Aalen--Johansen cumulative incidences \cite{aalen1978empirical} for recovery and lock-in, then a multi-state Cox model \cite{andersen1982cox,putter2007multistate} with position, mode, fault side, system and checkpoint distance as covariates, and a shared frailty per agent and per task \cite{vaupel1979frailty} for Q5. Validation on natural labels: $t_{\mathrm{obs}}-t_{\mathrm{err}}$ and $t_{\mathrm{lock}}-t_{\mathrm{err}}$ of the 1,184 Failure-as-a-Process failures, the detour length and recovery count of its 610 passes, and AgentRx's recovered events.
\item[E3 --- Checkpoint ablation.] Align $r(k/T)$ with $\mathcal{C}(\tau)$ in systems with an explicit verifier (Magentic-One's ledger, LongRCA's verification expert, the test runs of SWE agents), then remove, add or delay the verifier and re-run E1. The checkpoint model predicts the shift; the alternative---recovery driven by observations rather than by designated checks---predicts none.
\item[E4 --- Propagation matrix.] On the MASLab systems of Table~\ref{tab:engine} (row 4), inject at the output of agent $a$ and judge the wrongness of each consuming agent's next output against the unperturbed branch; estimate $M_{ab}$, $\rho(M)$ and the branching-process absorption probability, and compare with the measured $C$. The same data give $(g_a,\pi_a,c_a,x_a)$ per role.
\item[E5 --- Dose and interaction.] Severity-graded and paired injections; report $\sigma^\star(k)$ and the interaction term of Q7, on the same prefixes as E1.
\item[E6 --- Recovered-error detection from logs.] Train on outcome-only data (level L1 of Proposal~2) a hidden semi-Markov model \cite{yu2010hsmm} with an ``error live'' state that can switch back, and a value curve $V_t$ with a temporal-difference consistency loss, and score recovery segments against the four released sources of recovered errors: Failure-as-a-Process passes (first error step, recovery count), AgentRx non-root events (\S\ref{ds:agentrx}), TRAIL traces that end correctly but carry error spans (\S\ref{ds:trail}), and MP-Bench steps flagged by one annotator of three (\S\ref{ds:mp_bench}). Report recall of recovered errors separately from recall of fatal ones; the one-class detectors (\textsc{Oat}) are the baseline and are expected to flag both alike.
\item[E7 --- Online monitor with a delay guarantee.] A CUSUM \cite{page1954cusum} on the hazard or value sequence of U3, scored by lead time to $t_{\mathrm{lock}}$ and by false alarms on passes with recovered errors, the hardest negatives; the Lorden delay--false-alarm bound \cite{lorden1971} gives the curve a monitor cannot beat, against which the 82\% precision / 18\% recall / zero lead time of the Failure-as-a-Process prompted monitor can be placed.
\end{description}
\end{multicols}

\begin{table}[h]\centering\small
\begin{tabularx}{\linewidth}{@{}>{\RaggedRight\hsize=0.7\hsize}X >{\RaggedRight\hsize=0.95\hsize}X >{\RaggedRight\hsize=1.2\hsize}X >{\RaggedRight\hsize=1.15\hsize}X@{}}\toprule
Field & Model & Reading on agent trajectories & What it yields \\\midrule
Dependable computing \cite{avizienis2004dependability,arlat1990fault,voas1992pie} & fault / error / failure chain; coverage and latency; PIE & the fate taxonomy of Definition~\ref{def:fate} & robustness decomposed into execution, infection, propagation \\
Software reliability \cite{hiller2004epic} & error permeability and exposure per module & per-agent $(g_a,\pi_a,c_a,x_a)$, matrix $M_{ab}$ & which agent to harden, which edge needs a check \\
Event-history analysis \cite{putter2007multistate,aalen1978empirical,andersen1982cox,vaupel1979frailty} & competing risks, multi-state and recurrent-event models, frailty & states live / recovered / locked in; several errors per run & $F_{\mathrm{rec}}$, $F_{\mathrm{lock}}$, covariate effects, per-agent random effects, with censoring handled \\
Branching processes \cite{harris1963branching} & multi-type Galton--Watson, $\rho(M)$ & wrong message as an infection on the agent graph & analytic absorption probability; topology comparison \\
Ecology and infrastructure \cite{holling1973resilience,bruneau2003resilience} & basin of attraction; resilience triangle & $V_t$ as the state, lock-in as leaving the basin & critical severity $\sigma^\star$; depth and rapidity as separate numbers \\
Sequential analysis \cite{page1954cusum,lorden1971,tartakovsky2014sequential} & CUSUM, SPRT, quickest change detection & online monitor on the prefix (task U3) & detection delay against false-alarm rate, with a lower bound \\
Time series \cite{yu2010hsmm,adams2007bocpd} & hidden semi-Markov and switching models & latent on-track / derailed regime that can switch back & label-free recovery segments; Proposal~2's change point is the no-return case \\
Safe RL \cite{moldovan2012safe,grinsztajn2021reversibility} & recoverability and reversibility of states & point of no return as an irreversible state & learned lock-in predictors with a definition \\\bottomrule
\end{tabularx}
\caption{Mathematical models that the lifecycle programme borrows, and what each one contributes.}
\label{tab:lifecycle_models}
\end{table}

\begin{multicols}{2}
\paragraph{Relation to Proposals 1--3} The screen of Proposal~1 produces the rollouts; E1 adds the consumption test and keeps the branches, so the marginal cost over Proposal~1 is one LLM call per continuation step. Proposal~2's hazard model is the $h_{\mathrm{rec}}\equiv0$ case of Fig.~\ref{fig:lifecycle}; the multi-state version fits the PASS runs that Table~\ref{tab:contam} shows to be contaminated, and the recovered errors that Proposal~2 lists as a pitfall become a labelled target (E6). For Proposal~3, the fate labels change the credit: a recovered error should not carry the negative advantage of a fatal one, and the recovering step is a positive example of the capability the hypothesis names, which is what recovery-targeted training already exploits without localising it \cite{cuarecovery2026,reterm2026}.

\paragraph{Feasibility with the surveyed data}
\begin{itemize}[leftmargin=1.1em]
\item \emph{Released recovered errors} (for E2 validation and E6): Failure-as-a-Process ships, per PASS run, the first error step, the number of error events and recoveries, the maximal detour length and a near-miss flag, and per FAIL run $t_{\mathrm{err}}/t_{\mathrm{lock}}/t_{\mathrm{obs}}$, the number of recovery attempts, forewarning signals and the awareness class; AgentRx annotates every failure event including the recovered ones (up to 55 per trajectory); TRAIL's 148 runs are not selected by outcome; MP-Bench's three votes per step separate consensus errors from minority ones. No new annotation is needed for the natural side.
\item \emph{Injected side.} Aegis reports absorption rates but ships no absorbed branches, and Who\&When~Pro and TracerTraj discard them likewise, so E1 is a re-run of the released injectors (Aegis's factory, Who\&When~Pro's injector) with the negatives kept, on rows 1--5 of Table~\ref{tab:engine}; row 4 (MASLab systems, no environment state) is the cheapest and the natural home of E4.
\item \emph{Budget.} E1 reuses Proposal~1's $\approx$25\,k continuations and adds one call per continuation step ($\approx10^6$ calls, below the rollout cost); E2, E6 and E7 are statistical fits on the result; E3--E5 multiply the injection side by the number of ablations and severities (3--5$\times$) on a subset of prefixes.
\end{itemize}

\paragraph{Pitfalls to design for}
\begin{itemize}[leftmargin=1.1em]
\item \emph{Censoring.} A run that ends with the error live and $S{=}1$ carries no $R$ and no $L$; naive fractions over observed events are biased, which is why E2 uses survival estimators rather than averages. Short runs (ALFWorld's 30-step cap) censor heavily.
\item \emph{Dormancy inflates coverage.} An injection into content the task never needs is absorbed trivially; Definition~\ref{def:fate} separates it, and the dormant fraction per injector is itself a result about the injectors of Section~\ref{sec:g2}.
\item \emph{Detection is a judgement, recovery is a measurement.} $D$ depends on a judge reading the text; $R$ and $L$ are counterfactual. Report inter-judge agreement on $D$ and never let $D$ define the fate.
\item \emph{Baseline stochasticity.} 39 of 54 successful MAS runs fail at least once in three reruns \cite{luan2026symtrace}, so $V^0$ is a curve with intervals, and ``fatal'' is a drop relative to it, not an absolute failure.
\item \emph{Wrongness of downstream outputs} (E4) needs a reference: the unperturbed branch is stochastic, so the judge must compare against the ground truth or a sample of unperturbed continuations, with the same caveat as $D$.
\item \emph{Injector realism and position prior.} The released injectors place errors late (Fig.~\ref{fig:errpos}); inject uniformly over positions and stratify every curve by $k/T$, or the lifetime curves inherit the prior.
\item \emph{Symptom shift in natural labels.} $t_{\mathrm{obs}}$ can precede $t_{\mathrm{lock}}$ in Failure-as-a-Process and judge labels drift to the symptom \cite{shen2026epv}; the natural data validate the \emph{shape} of the curves, not their location.
\end{itemize}

\paragraph{What would count as a result}
\begin{itemize}[leftmargin=1.1em]
\item The first measured lifetime curves of errors in agent systems: recovery and lock-in incidences by lag, position, mode and system, with the fraction of absorption that is genuine correction rather than task slack.
\item A tested checkpoint model: either $r(k/T)$ is a step function that a verifier ablation shifts by the predicted amount, or recovery is observation-driven and verifier agents are not where robustness comes from.
\item Per-agent profiles $(g_a,\pi_a,c_a,x_a)$ for the roles of the surveyed systems and a propagation matrix whose spectral radius predicts absorption, which would turn ``which agent is responsible'' into ``which agent to harden''.
\item A label-free detector of recovered errors with a measured delay--false-alarm curve for the online monitor, placed against the Lorden bound.
\end{itemize}
\end{multicols}
